\section{Discussion and Conclusion}
\label{sec:discussion}
\subsection{Lessons for Link-State Forecasting}
Four lessons apply beyond our testbed. First, at the sub-second scale the level of LEO and 5G RTT is close to unpredictable beyond a return toward its recent value, while its spread is predictable from the recent input. Forecasters should be judged by what is predictable. A point metric such as the RMSE hardly separates a recurrent network from a linear model or from its own predecessor, while a calibrated interval score does. Second, the calibration must use sessions that the model has not seen, and the coverage must be checked per volatility regime. Intervals from the test residuals hide the difference between forecasters, and marginal coverage hides the misallocation of an interval of constant width. Third, measurement pipelines need a causality check. A time interpolation that is harmless for plots places future samples in the inputs, and its effect is largest at the shortest horizons, where schedulers operate. Fourth, a quantile forecaster with conformal calibration across sessions is the tool that captures what is predictable. Its linear version is a strong reference that any deep forecaster should be compared against~\cite{zeng2023transformers}.

\subsection{Use in a Multipath Scheduler}
The calibrated upper quantile is a per-path RTT bound that a deadline-aware scheduler can use directly. It sends a packet on one path if that path's bound at the delivery horizon meets the deadline and duplicates it otherwise. The bound is tight on a quiet path and loose only when the path is volatile. The median orders the paths by their expected delay, and the mean suits capacity estimates, where the squared error is the right loss. All three come from one forward pass. The conformal guarantee of a bound holds on average and only if calibration and test sessions are exchangeable. On our test sessions the one-sided 95\% bounds covered \BoundCovLo{} to \BoundCovHi{} of the RTT values. Adaptive conformal updates~\cite{gibbs2021adaptive} can correct such drift online.

\subsection{Limitations}
The sessions were recorded at one site with a fixed terminal, one ping server, and constant-rate iperf traffic, and the conditions of a moving or airborne terminal may differ. The OneWeb sessions were recorded in 2025 and the Starlink and 5G sessions in 2026. The scheduling study therefore pairs paths that were not measured at the same time and assumes that they are independent. Starlink UL shows a shift between calibration and test sessions that no method corrects. DS2S-E under-covers the rare OneWeb bursts, and the anchored decoder ranks OneWeb DL spikes worse than the previous model. The scheduling benefit is emulated from measured traces, and an end-to-end evaluation inside the MPQUIC client is the next step.

\subsection{Conclusion}
We asked what can be forecast about the sub-second RTT of Starlink, OneWeb, and 5G links. The next change is close to unpredictable, but the spread of the RTT depends on the link state, which the recent input reveals, and its errors are skewed upward. A calibrated distributional forecaster turns this into intervals that are narrow on quiet links and wide on volatile ones. Against the same network with a scalar output, it improves the 95\% interval score on all six links by \BSISWorstAbs\% to \BSISBestAbs\%, significantly at the session level, and it narrows the coverage gap between calm and volatile windows on four links. A linear quantile model recovers most of this gain, which places the value in the calibrated distributional output. In an emulated deadline-aware scheduler, the bound reduces duplication on Starlink DL by \SchedSaveSLDL\% without more late packets. Future work will close the loop with the multipath scheduler, add the state of the other paths as inputs, and make the calibration aware of the regime.
